% Part: model-theory
% Chapter: basics
% Section: partial-iso

\documentclass[../../../include/open-logic-section]{subfiles}

\begin{document}

\olfileid{mod}{bas}{pis}
\section{التشاكلات الجزئية}

\begin{defn}
  لتكن !!{structure}s اثنتان $\Struct{M}$ و$\Struct{N}$. \emph{التشاكل
  الجزئي} من $\Struct{M}$ إلى $\Struct{N}$ دالة جزئية منتهية~${\OLId{latin-ordinary}{p}}$ تأخذ
  مدخلاتها من $\Domain {\OLId{latin-looped}{M}}$ وتعيد قيمًا في $\Domain {\OLId{latin-looped}{N}}$، وتحقق شروط التشاكل
  الواردة في \olref[iso]{defn:isomorphism} على مجالها:
  \begin{enumerate}
  \item تكون ${\OLId{latin-ordinary}{p}}$~!!{injective}؛
  \item لكل !!{constant}~${\OLId{latin-ordinary}{c}}$: إذا كانت ${\OLId{latin-ordinary}{p}}(\Assign{{\OLId{latin-ordinary}{c}}}{{\OLId{latin-looped}{M}}})$ معرّفة، فإن
    ${\OLId{latin-ordinary}{p}}(\Assign{{\OLId{latin-ordinary}{c}}}{{\OLId{latin-looped}{M}}}) = \Assign{{\OLId{latin-ordinary}{c}}}{{\OLId{latin-looped}{N}}}$؛
  \item لكل !!{predicate}~${\OLId{latin-looped}{P}}$ ذي ${\OLId{latin-ordinary}{n}}$ مواضع: إذا كانت ${\OLId{latin-ordinary}{a}}_{\OLMathNumeral{1}}$، \dots،
    ${\OLId{latin-ordinary}{a}}_{\OLId{latin-ordinary}{n}}$ في مجال ${\OLId{latin-ordinary}{p}}$، فإن $\langle {\OLId{latin-ordinary}{a}}_{\OLMathNumeral{1}}, \dots, {\OLId{latin-ordinary}{a}}_{\OLId{latin-ordinary}{n}}\rangle \in
    \Assign {\OLId{latin-looped}{P}} {\OLId{latin-looped}{M}}$ إذا وفقط إذا كان
    $\langle {\OLId{latin-ordinary}{p}}({\OLId{latin-ordinary}{a}}_{\OLMathNumeral{1}}), \dots, {\OLId{latin-ordinary}{p}}({\OLId{latin-ordinary}{a}}_{\OLId{latin-ordinary}{n}}) \rangle \in \Assign {\OLId{latin-looped}{P}} {\OLId{latin-looped}{N}}$؛
  \item لكل !!{function}~${\OLId{latin-ordinary}{f}}$ ذات ${\OLId{latin-ordinary}{n}}$ مواضع: إذا كانت ${\OLId{latin-ordinary}{a}}_{\OLMathNumeral{1}}$، \dots،
    ${\OLId{latin-ordinary}{a}}_{\OLId{latin-ordinary}{n}}$ في مجال ${\OLId{latin-ordinary}{p}}$، فإن ${\OLId{latin-ordinary}{p}}(\Assign {\OLId{latin-ordinary}{f}} {\OLId{latin-looped}{M}} ({\OLId{latin-ordinary}{a}}_{\OLMathNumeral{1}}, \dots,{\OLId{latin-ordinary}{a}}_{\OLId{latin-ordinary}{n}}))
    = \Assign {\OLId{latin-ordinary}{f}} {\OLId{latin-looped}{N}} ({\OLId{latin-ordinary}{p}}({\OLId{latin-ordinary}{a}}_{\OLMathNumeral{1}}), \dots, {\OLId{latin-ordinary}{p}}({\OLId{latin-ordinary}{a}}_{\OLId{latin-ordinary}{n}}))$.
  \end{enumerate}
  وكون ${\OLId{latin-ordinary}{p}}$ منتهية يعني أن $\dom{{\OLId{latin-ordinary}{p}}}$ منتهٍ.
\end{defn}

ومن ذلك أن الدالة الخالية~$\emptyset$ تعدّ تشاكلًا جزئيًا بين كل
!!{structure}s اثنتين.

\begin{defn}\ollabel{defn:partialisom}
  تكون !!{structure}s اثنتان $\Struct{M}$ و$\Struct{N}$
  \emph{متشاكلتين جزئيًا}، ونكتب $\Struct{M} \iso[{\OLId{latin-ordinary}{p}}] \Struct{N}$، إذا
  وفقط إذا وجدت مجموعة غير خالية~${\OLId{latin-looped}{I}}$ من التشاكلات الجزئية بين
  $\Struct{M}$ و$\Struct{N}$ تحقق خاصية \emph{الذهاب والإياب}:
  \begin{enumerate}
  \item (\emph{الذهاب}) لكل ${\OLId{latin-ordinary}{p}} \in {\OLId{latin-looped}{I}}$ و${\OLId{latin-ordinary}{a}} \in \Domain {\OLId{latin-looped}{M}}$ يوجد
    ${\OLId{latin-ordinary}{q}} \in {\OLId{latin-looped}{I}}$ بحيث ${\OLId{latin-ordinary}{p}} \subseteq {\OLId{latin-ordinary}{q}}$ ويكون ${\OLId{latin-ordinary}{a}}$ في مجال ${\OLId{latin-ordinary}{q}}$؛
  \item (\emph{الإياب}) لكل ${\OLId{latin-ordinary}{p}} \in {\OLId{latin-looped}{I}}$ و${\OLId{latin-ordinary}{b}} \in \Domain {\OLId{latin-looped}{N}}$ يوجد
    ${\OLId{latin-ordinary}{q}} \in {\OLId{latin-looped}{I}}$ بحيث ${\OLId{latin-ordinary}{p}} \subseteq {\OLId{latin-ordinary}{q}}$ ويكون ${\OLId{latin-ordinary}{b}}$ في مدى ${\OLId{latin-ordinary}{q}}$.
  \end{enumerate}
\end{defn}

\begin{thm}\ollabel{thm:p-isom1}
  إذا كان $\Struct{M} \iso[{\OLId{latin-ordinary}{p}}] \Struct{N}$ وكانت $\Struct{M}$ و
  $\Struct{N}$~!!{enumerable}، فإن $\Struct{M} \iso \Struct{N}$.
\end{thm}

\begin{proof}
  بما أن $\Struct{M}$ و$\Struct{N}$~!!{enumerable}، فلنكتب
  $\Domain{{\OLId{latin-looped}{M}}} = \{{\OLId{latin-ordinary}{a}}_{\OLMathNumeral{0}}, {\OLId{latin-ordinary}{a}}_{\OLMathNumeral{1}}, \ldots \}$ و
  $\Domain{{\OLId{latin-looped}{N}}} = \{{\OLId{latin-ordinary}{b}}_{\OLMathNumeral{0}}, {\OLId{latin-ordinary}{b}}_{\OLMathNumeral{1}}, \ldots \}$. نختار أولًا ${\OLId{latin-ordinary}{p}}_{\OLMathNumeral{0}} \in {\OLId{latin-looped}{I}}$،
  ثم نبني تشاكلات جزئية متزايدة
  ${\OLId{latin-ordinary}{p}}_{\OLMathNumeral{0}} \subseteq {\OLId{latin-ordinary}{p}}_{\OLMathNumeral{1}} \subseteq {\OLId{latin-ordinary}{p}}_{\OLMathNumeral{2}} \subseteq \cdots$، فنناوب
  بين الذهاب والإياب على الوجه الآتي:
  \begin{enumerate}
  \item إذا كان ${\OLId{latin-ordinary}{n}}+{\OLMathNumeral{1}}$ فرديًا، أي إذا كان ${\OLId{latin-ordinary}{n}} = {\OLMathNumeral{2}}{\OLId{latin-ordinary}{r}}$، فاستعمل خاصية الذهاب
    لإيجاد ${\OLId{latin-ordinary}{p}}_{{\OLId{latin-ordinary}{n}}+{\OLMathNumeral{1}}} \in {\OLId{latin-looped}{I}}$ بحيث ${\OLId{latin-ordinary}{p}}_{\OLId{latin-ordinary}{n}} \subseteq {\OLId{latin-ordinary}{p}}_{{\OLId{latin-ordinary}{n}}+{\OLMathNumeral{1}}}$ ويكون ${\OLId{latin-ordinary}{a}}_{\OLId{latin-ordinary}{r}}$
    في مجال ${\OLId{latin-ordinary}{p}}_{{\OLId{latin-ordinary}{n}}+{\OLMathNumeral{1}}}$؛
  \item إذا كان ${\OLId{latin-ordinary}{n}}+{\OLMathNumeral{1}}$ زوجيًا، أي إذا كان ${\OLId{latin-ordinary}{n}}={\OLMathNumeral{2}}{\OLId{latin-ordinary}{r}}+{\OLMathNumeral{1}}$، فاستعمل خاصية الإياب
    لإيجاد ${\OLId{latin-ordinary}{p}}_{{\OLId{latin-ordinary}{n}}+{\OLMathNumeral{1}}} \in {\OLId{latin-looped}{I}}$ بحيث ${\OLId{latin-ordinary}{p}}_{\OLId{latin-ordinary}{n}} \subseteq {\OLId{latin-ordinary}{p}}_{{\OLId{latin-ordinary}{n}}+{\OLMathNumeral{1}}}$ ويكون ${\OLId{latin-ordinary}{b}}_{\OLId{latin-ordinary}{r}}$
    في مدى ${\OLId{latin-ordinary}{p}}_{{\OLId{latin-ordinary}{n}}+{\OLMathNumeral{1}}}$.
  \end{enumerate}
ثم نضع:
\[
{\OLId{latin-ordinary}{p}} = \bigcup_{{\OLId{latin-ordinary}{n}}\ge {\OLMathNumeral{0}}} {\OLId{latin-ordinary}{p}}_{\OLId{latin-ordinary}{n}},
\]
كان ${\OLId{latin-ordinary}{p}}$ تشاكلًا بين $\Struct{M}$ و$\Struct{N}$.
\end{proof}

\begin{prob}
  استوفِ تفاصيل البرهان على أن ${\OLId{latin-ordinary}{p}}$ المعرفة في
  \olref[mod][bas][pis]{thm:p-isom1} تشاكل.
\end{prob}

\begin{thm}\ollabel{thm:p-isom2}
  لتكن $\Struct{M}$ و$\Struct{N}$~!!{structure}s للـ!!{language} علائقية
  محضة (أي !!a{language} لا تحتوي إلا !!{predicate}s، ولا تحتوي
  !!{function}s أو ثوابت). إذا كان $\Struct{M} \iso[{\OLId{latin-ordinary}{p}}] \Struct{N}$،
  فإن $\Struct{M} \elemequiv \Struct{N}$ أيضًا.
\end{thm}

\begin{proof}
  % تصحيح مصدر (0189-I1): قُيّد التشاكل الجزئي بعائلة الذهاب والإياب، وقُيّدت المتغيرات الحرة، وصُرّح بحجّة المدّ التي أغفلها الأصل وخط الأساس.
  اختر عائلةً ${\OLId{latin-looped}{I}}$ من التشاكلات الجزئية تشهد خاصية الذهاب والإياب.
  ونثبت بالاستقراء على الـ!!{formula}s الدعوى الآتية: إذا كانت
  متغيرات~$!A$ الحرة محصورة في ${\OLId{latin-ordinary}{x}}_{\OLMathNumeral{1}}$، \dots،~${\OLId{latin-ordinary}{x}}_{\OLId{latin-ordinary}{n}}$، وكان ${\OLId{latin-ordinary}{p}}\in {\OLId{latin-looped}{I}}$
  يرسل كل ${\OLId{latin-ordinary}{a}}_{\OLId{latin-ordinary}{i}}$ إلى ${\OLId{latin-ordinary}{b}}_{\OLId{latin-ordinary}{i}}$، وكان ${\OLId{latin-ordinary}{s}}_{\OLMathNumeral{1}}({\OLId{latin-ordinary}{x}}_{\OLId{latin-ordinary}{i}})={\OLId{latin-ordinary}{a}}_{\OLId{latin-ordinary}{i}}$ و${\OLId{latin-ordinary}{s}}_{\OLMathNumeral{2}}({\OLId{latin-ordinary}{x}}_{\OLId{latin-ordinary}{i}})={\OLId{latin-ordinary}{b}}_{\OLId{latin-ordinary}{i}}$
  لكل ${\OLId{latin-ordinary}{i}}={\OLMathNumeral{1}}$، \dots،~${\OLId{latin-ordinary}{n}}$، فإن $\Sat{{\OLId{latin-looped}{M}}}{!A}[{\OLId{latin-ordinary}{s}}_{\OLMathNumeral{1}}]$ إذا وفقط إذا كان
  $\Sat{{\OLId{latin-looped}{N}}}{!A}[{\OLId{latin-ordinary}{s}}_{\OLMathNumeral{2}}]$. الحالة الذرية من كون~${\OLId{latin-ordinary}{p}}$ تشاكلًا جزئيًا،
  وحالات الروابط مباشرة. وفي حالتي المكمّمين نستعمل خاصيتي الذهاب
  والإياب لمدّ~${\OLId{latin-ordinary}{p}}$ إلى ${\OLId{latin-ordinary}{q}}\in {\OLId{latin-looped}{I}}$ يضم الشاهد الجديد في مجاله أو مداه،
  ثم نطبق فرض الاستقراء. فإذا كان~$!A$ جملة أخذنا ${\OLId{latin-ordinary}{n}}={\OLMathNumeral{0}}$ وأي~${\OLId{latin-ordinary}{p}}\in {\OLId{latin-looped}{I}}$،
  فحصل $\Struct{M} \elemequiv \Struct{N}$.
\end{proof}

\begin{rem}
ولا تبطل النتيجة إذا كانت في اللغة !!{function}s؛ وإنما ننظر
حينئذ في التشاكل الذي تولّده ${\OLId{latin-ordinary}{p}}$ بين الـ!!{structure} الجزئية
من $\Struct{M}$ المولّدة بالعناصر ${\OLId{latin-ordinary}{a}}_{\OLMathNumeral{1}}$، \dots، ${\OLId{latin-ordinary}{a}}_{\OLId{latin-ordinary}{n}}$،
والـ!!{structure} الجزئية من $\Struct{N}$ المولّدة بالعناصر
${\OLId{latin-ordinary}{b}}_{\OLMathNumeral{1}}$، \dots، ${\OLId{latin-ordinary}{b}}_{\OLId{latin-ordinary}{n}}$.
\end{rem}

ولنا أن نفصّل النتيجة السابقة مرحلةً مرحلة: نقابل بين عدد
المكممات المتداخلة في !!a{formula} وعدد المرات التي نستطيع فيها
توسيع التشاكلات الجزئية. ولأجل ذلك نضع التعريفات الآتية.

\begin{defn}
  إذا أعطيت !!{formula}~$!A$، فـ\emph{رتبة الكم} لـ$!A$، ورمزها
  $\QuantRank{!A} \in \Nat$، تحدّ بالاستقراء بأكبر عدد من المكممات
  المتداخلة في $!A$. وتكون !!{structure}s اثنتان $\Struct{M}$ و
  $\Struct{N}$ \emph{متكافئتين حتى الرتبة~${\OLId{latin-ordinary}{n}}$}، ونكتب
  $\Struct{M} \elemequiv[{\OLId{latin-ordinary}{n}}] \Struct{N}$، إذا اتفقتا على جميع الجمل ذات
  رتبة الكم الأصغر من أو المساوية لـ${\OLId{latin-ordinary}{n}}$.
\end{defn}

\begin{prop}\ollabel{prop:qr-finite}
  لتكن $\Lang{L}$~!!{language} علائقية محضة منتهية، أي !!{language}
  تحتوي عددًا منتهيًا من الـ!!{predicate}s والـ!!{constant}s ولا تحتوي
  !!{function}s. لكل ${\OLId{latin-ordinary}{n}},{\OLId{latin-ordinary}{k}} \in \Nat$ لا يوجد، حتى التكافؤ المنطقي، إلا
  عدد منتهٍ من !!{formula}s الرتبة الأولى في الـ!!{language}~$\Lang{L}$ التي لا تزيد
  رتبة كمها على~${\OLId{latin-ordinary}{n}}$ وتقع متغيراتها الحرة ضمن ${\OLId{latin-ordinary}{x}}_{\OLMathNumeral{1}}$، \dots،~${\OLId{latin-ordinary}{x}}_{\OLId{latin-ordinary}{k}}$.
\end{prop}

\begin{proof}
  نستقرئ على ${\OLId{latin-ordinary}{n}}$، ونثبت الدعوى في كل مرحلة لجميع قيم ${\OLId{latin-ordinary}{k}}$ معًا.
\end{proof}

\begin{defn}
  لتكن !!a{structure}~$\Struct{M}$. نرمز بـ$\Domain {\OLId{latin-looped}{M}}^{<\omega}$ إلى
  مجموعة جميع المتتاليات المنتهية على $\Domain{{\OLId{latin-looped}{M}}}$. ونستعمل
  $\mathbf{a}, \mathbf{b}, \mathbf{c}, \ldots$ للدلالة على متتاليات
  منتهية. وإذا كانت $\mathbf{a} \in \Domain{{\OLId{latin-looped}{M}}}^{<\omega}$ و
  ${\OLId{latin-ordinary}{a}} \in \Domain{{\OLId{latin-looped}{M}}}$، فإن $\mathbf{a}{\OLId{latin-ordinary}{a}}$ تمثل \emph{وصل}
  $\mathbf{a}$ بـ${\OLId{latin-ordinary}{a}}$.
\end{defn}

\begin{defn}
  لتكن !!{structure}s اثنتان $\Struct{M}$ و$\Struct{N}$. نعرّف علائق
  ${\OLId{latin-looped}{I}}_{\OLId{latin-ordinary}{n}} \subseteq \Domain {\OLId{latin-looped}{M}}^{<\omega} \times \Domain {\OLId{latin-looped}{N}}^{<\omega}$ بين
  متتاليات متساوية الطول، بالاستقراء على ${\OLId{latin-ordinary}{n}}$ كما يأتي:
   \begin{enumerate}
   \item تكون ${\OLId{latin-looped}{I}}_{\OLMathNumeral{0}}(\mathbf{a},\mathbf{b})$ إذا وفقط إذا حققت
     $\mathbf{a}$ و$\mathbf{b}$ الـ!!{formula}s الذرية نفسها في
     $\Struct{M}$ و$\Struct{N}$؛ أي إذا كان طول المتتاليتين~${\OLId{latin-ordinary}{k}}$، وكان
     ${\OLId{latin-ordinary}{s}}_{\OLMathNumeral{1}}({\OLId{latin-ordinary}{x}}_{\OLId{latin-ordinary}{i}})={\OLId{latin-ordinary}{a}}_{\OLId{latin-ordinary}{i}}$ و${\OLId{latin-ordinary}{s}}_{\OLMathNumeral{2}}({\OLId{latin-ordinary}{x}}_{\OLId{latin-ordinary}{i}})={\OLId{latin-ordinary}{b}}_{\OLId{latin-ordinary}{i}}$، وكانت $!A$ ذرية وجميع
     !!{variable}s فيها من بين ${\OLId{latin-ordinary}{x}}_{\OLMathNumeral{1}}$، \dots،~${\OLId{latin-ordinary}{x}}_{\OLId{latin-ordinary}{k}}$، فإن
     $\Sat{{\OLId{latin-looped}{M}}}{!A}[{\OLId{latin-ordinary}{s}}_{\OLMathNumeral{1}}]$ إذا وفقط إذا كان $\Sat{{\OLId{latin-looped}{N}}}{!A}[{\OLId{latin-ordinary}{s}}_{\OLMathNumeral{2}}]$.
   \item تكون ${\OLId{latin-looped}{I}}_{{\OLId{latin-ordinary}{n}}+{\OLMathNumeral{1}}}(\mathbf{a},\mathbf{b})$ إذا وفقط إذا كان لكل
     ${\OLId{latin-ordinary}{a}}\in \Domain {\OLId{latin-looped}{M}}$ عنصر~${\OLId{latin-ordinary}{b}}\in \Domain {\OLId{latin-looped}{N}}$ بحيث
     ${\OLId{latin-looped}{I}}_{\OLId{latin-ordinary}{n}}(\mathbf{a}{\OLId{latin-ordinary}{a}},\mathbf{b}{\OLId{latin-ordinary}{b}})$، والعكس بالعكس.
   \end{enumerate}
\end{defn}


\begin{defn}
  نكتب $\Struct{M} \approx_{\OLId{latin-ordinary}{n}} \Struct{N}$ إذا تحققت
  ${\OLId{latin-looped}{I}}_{\OLId{latin-ordinary}{n}}(\emptyseq,\emptyseq)$ في $\Struct{M}$ و$\Struct{N}$ (حيث
  $\emptyseq$ هي المتتالية الخالية).
\end{defn}

\begin{thm}\ollabel{thm:b-n-f}
  لتكن $\Lang{L}$~!!{language} علائقية محضة. إذا كانت
  ${\OLId{latin-looped}{I}}_{\OLId{latin-ordinary}{n}}(\mathbf{a},\mathbf{b})$ وكان طول المتتاليتين~${\OLId{latin-ordinary}{k}}$، فإن لكل $!A$
  تقع متغيراتها الحرة ضمن ${\OLId{latin-ordinary}{x}}_{\OLMathNumeral{1}}$، \dots،~${\OLId{latin-ordinary}{x}}_{\OLId{latin-ordinary}{k}}$ وتحقق
  $\QuantRank{!A} \le {\OLId{latin-ordinary}{n}}$، يكون $\Sat{{\OLId{latin-looped}{M}}}{!A}[\mathbf{a}]$ إذا وفقط إذا
  كان $\Sat{{\OLId{latin-looped}{N}}}{!A}[\mathbf{b}]$ (حيث تحقق $\mathbf{a}$ مرة أخرى~$!A$
  إذا حققها أي تعيين~${\OLId{latin-ordinary}{s}}$ بحيث ${\OLId{latin-ordinary}{s}}({\OLId{latin-ordinary}{x}}_{\OLId{latin-ordinary}{i}})={\OLId{latin-ordinary}{a}}_{\OLId{latin-ordinary}{i}}$). وإذا كانت $\Lang{L}$
  منتهية، صح العكس أيضًا.
\end{thm}

\begin{proof}
  أما جهة اللزوم، فاستقراء يسير على $!A$ يبيّن أن
  ${\OLId{latin-looped}{I}}_{\OLId{latin-ordinary}{n}}(\mathbf{a},\mathbf{b})$ تستلزم اتفاق $\mathbf{a}$ و$\mathbf{b}$
  في الـ!!{formula}s التي لا تزيد رتبة كمها على~${\OLId{latin-ordinary}{n}}$، والتي تقع
  متغيراتها الحرة في القائمة الثابتة الموافقة لإحداثياتهما.
  وأما العكس، فنستقرئ فيه على ${\OLId{latin-ordinary}{n}}$، ونعتمد
  \olref{prop:qr-finite}. فهي تضمن، لكل ${\OLId{latin-ordinary}{n}}$ ولكل قائمة منتهية
  ثابتة من المتغيرات الحرة، أن فئات التكافؤ المنطقي
  للـ!!{formula}s ذات رتبة الكم التي لا تتجاوز~${\OLId{latin-ordinary}{n}}$ منتهية العدد.

  عند ${\OLId{latin-ordinary}{n}}={\OLMathNumeral{0}}$، نعلم بالفرض أن $\mathbf{a}$ و$\mathbf{b}$ تتفقان
  في الـ!!{formula}s الخالية من الكم؛ فتتفقان خاصة في الصيغ
  الذرية، وهذا بعينه ${\OLId{latin-looped}{I}}_{\OLMathNumeral{0}}(\mathbf{a},\mathbf{b})$.

  وللحالة ${\OLId{latin-ordinary}{n}}+{\OLMathNumeral{1}}$ نفرض اتفاق $\mathbf{a}$ و$\mathbf{b}$ في
  الـ!!{formula}s التي لا تزيد رتبة كمها على ${\OLId{latin-ordinary}{n}}+{\OLMathNumeral{1}}$.
  لإثبات ${\OLId{latin-looped}{I}}_{{\OLId{latin-ordinary}{n}}+{\OLMathNumeral{1}}}(\mathbf{a},\mathbf{b})$ ننظر أولًا في جهة
  الذهاب: لكل ${\OLId{latin-ordinary}{a}} \in \Domain {\OLId{latin-looped}{M}}$ نطلب ${\OLId{latin-ordinary}{b}} \in \Domain {\OLId{latin-looped}{N}}$ بحيث
  ${\OLId{latin-looped}{I}}_{\OLId{latin-ordinary}{n}}(\mathbf{a}{\OLId{latin-ordinary}{a}},\mathbf{b}{\OLId{latin-ordinary}{b}})$. وبفرض الاستقراء يكفي أن
  نجد، لكل ${\OLId{latin-ordinary}{a}} \in \Domain {\OLId{latin-looped}{M}}$، عنصرًا ${\OLId{latin-ordinary}{b}} \in \Domain {\OLId{latin-looped}{N}}$
  تتفق عنده $\mathbf{a}{\OLId{latin-ordinary}{a}}$ و$\mathbf{b}{\OLId{latin-ordinary}{b}}$ في الـ!!{formula}s
  التي لا تزيد رتبة كمها على~${\OLId{latin-ordinary}{n}}$.

  فلنأخذ ${\OLId{latin-ordinary}{a}} \in \Domain {\OLId{latin-looped}{M}}$. نختار ممثلًا لكل فئة تكافؤ منطقي
  للـ!!{formula}s~$!B({\OLId{latin-ordinary}{x}},\mathbf{y})$ التي لا تزيد رتبة كمها على~${\OLId{latin-ordinary}{n}}$،
  وتقع متغيراتها الحرة في القائمة الثابتة ${\OLId{latin-ordinary}{x}},\mathbf{y}$،
  وتحققها $\mathbf{a}{\OLId{latin-ordinary}{a}}$ في $\Struct{M}$. ونجعل $!T^{\OLId{latin-ordinary}{a}}_{\OLId{latin-ordinary}{n}}$
  اقتران الممثلين المختارين. وهم عدد منتهٍ، فيكون الاقتران
  !!{formula} واحدة من الرتبة الأولى. وعندئذ تحقق $\mathbf{a}$
  الصيغة $\lexists[{\OLId{latin-ordinary}{x}}][!T^{\OLId{latin-ordinary}{a}}_{\OLId{latin-ordinary}{n}}({\OLId{latin-ordinary}{x}},\mathbf{y})]$، ولا تزيد رتبة
  كمها على ${\OLId{latin-ordinary}{n}}+{\OLMathNumeral{1}}$. فيلزم بالفرض أن تحقق $\mathbf{b}$ الصيغة
  نفسها في $\Struct{N}$؛ فيوجد ${\OLId{latin-ordinary}{b}} \in \Domain {\OLId{latin-looped}{N}}$ تحقق عنده
  $\mathbf{b}{\OLId{latin-ordinary}{b}}$ الصيغة $!T^{\OLId{latin-ordinary}{a}}_{\OLId{latin-ordinary}{n}}$. ومن ثم تتفق $\mathbf{b}{\OLId{latin-ordinary}{b}}$
  في الـ!!{formula}s ذات رتبة الكم التي لا تتجاوز~${\OLId{latin-ordinary}{n}}$
  مع $\mathbf{a}{\OLId{latin-ordinary}{a}}$؛ إذ في المختار ممثل لكل صيغة صادقة،
  أو لنفيها إن كانت كاذبة. وبنظير هذا النظر نثبت جهة الإياب:
  لكل ${\OLId{latin-ordinary}{b}} \in \Domain {\OLId{latin-looped}{N}}$ يوجد ${\OLId{latin-ordinary}{a}}\in \Domain {\OLId{latin-looped}{M}}$ تتفق عنده
  $\mathbf{a}{\OLId{latin-ordinary}{a}}$ و$\mathbf{b}{\OLId{latin-ordinary}{b}}$ في الـ!!{formula}s ذات رتبة
  الكم التي لا تتجاوز~${\OLId{latin-ordinary}{n}}$. فتم البرهان.
\end{proof}

\begin{cor}\ollabel{cor:b-n-f}
  إذا كانت $\Struct{M}$ و$\Struct{N}$~!!{structure}s علائقية محضة في
  !!{language} منتهية، فإن $\Struct{M} \approx_{\OLId{latin-ordinary}{n}}\Struct{N}$ إذا وفقط
  إذا كان $\Struct{M} \elemequiv[{\OLId{latin-ordinary}{n}}] \Struct{N}$. وبوجه خاص،
  $\Struct{M} \elemequiv \Struct{N}$ إذا وفقط إذا كان، لكل ${\OLId{latin-ordinary}{n}}$،
  $\Struct{M} \approx_{\OLId{latin-ordinary}{n}} \Struct{N}$.
\end{cor}

\end{document}
